% Generated draft: requires mathematical and source review.
\documentclass[11pt]{article}
\usepackage[margin=25mm]{geometry}
\usepackage{fontspec}
\setmainfont{Latin Modern Roman}
\usepackage{amsmath,amssymb,mathtools,unicode-math}
\setmathfont{Latin Modern Math}
\usepackage{enumitem,needspace,xurl,hyperref}
\setlength{\parindent}{0pt}
\setlength{\parskip}{5pt}
\setlength{\emergencystretch}{3em}
\newcommand{\N}{\mathbb{N}}
\newcommand{\Z}{\mathbb{Z}}
\newcommand{\Q}{\mathbb{Q}}
\newcommand{\R}{\mathbb{R}}
\newcommand{\C}{\mathbb{C}}
\newcommand{\F}{\mathbb{F}}
\newcommand{\A}{\mathbb{A}}
\newcommand{\PP}{\mathbb{P}}
\newcommand{\E}{\mathbb{E}}
\newcommand{\eps}{\varepsilon}
\newcommand{\dd}{\,\mathrm{d}}
\begin{document}
\section*{UnsolvedMath ID: 700686}
\textbf{Source problem code:} catalogue entry
\subsection{Definitions and mathematical statement}
Let \(F\) be a free group of finite rank, and let \(\operatorname{Aut}(F)\) be its group of automorphisms under composition. For a finite subset \(S\subseteq\operatorname{Aut}(F)\), define its centralizer by
\[
C_{\operatorname{Aut}(F)}(S)=\{\beta\in\operatorname{Aut}(F):\beta\circ\alpha=\alpha\circ\beta\text{ for every }\alpha\in S\}.
\]
For a single automorphism \(\alpha\), write \(C_{\operatorname{Aut}(F)}(\alpha)=C_{\operatorname{Aut}(F)}(\{\alpha\})\). A group is finitely presented if it has a presentation with finitely many generators and finitely many defining relations.

The source asks the following two questions:
\begin{enumerate}
\item For every free group \(F\) of finite rank and every \(\alpha\in\operatorname{Aut}(F)\), is \(C_{\operatorname{Aut}(F)}(\alpha)\) finitely presented?
\item More generally, for every free group \(F\) of finite rank and every finite subset \(S\subseteq\operatorname{Aut}(F)\), is \(C_{\operatorname{Aut}(F)}(S)\) finitely presented?
\end{enumerate}
An inner automorphism of \(F\) is an automorphism of the form \(x\mapsto gxg^{-1}\) for some \(g\in F\). The source reports an affirmative answer in rank \(2\) and in the case of inner automorphisms for every finite rank. Its affirmative-case remark does not separately discuss the single-automorphism and finite-set formulations; this scope qualification remains to be checked against the cited source. These reported cases are supplied background, not independently verified assertions in this draft.
\subsection{Short English statement}
For a free group of finite rank, must the group of automorphisms commuting with a specified automorphism be finitely presented? The source also asks the stronger question for automorphisms commuting with every member of a specified finite set. It reports affirmative cases in rank \(2\) and for inner automorphisms.
\subsection{Sources}
\begin{itemize}
\item Supplied local UnsolvedMath registry record: permanent ID \(700686\), exact source problem code \texttt{catalogue entry}, title \emph{Kourovka Notebook Problem 10.71}. The registry labels the problem \texttt{open}; this is supplied classification, not an independently verified literature assessment.
\item \emph{The Kourovka Notebook (Unsolved Problems in Group Theory)}, original problem 10.71, 10th issue (1986); proposed by V. A. Churkin. The supplied background identifies the 21st issue edition (2026), \texttt{arXiv:1401.0300v48}. The supplied collection description names E. I. Khukhro and V. D. Mazurov as editors.
\item The registry record was created and updated on \texttt{2026-10-16T14:00:00Z}. These are registry timestamps, not dates of independent source verification. No separate literature-check date or literature-source links are supplied.
\end{itemize}
\end{document}
